\section{Boundaries and Future Routes}

The grammar leaves several routes outside the current claim boundary.  The signed-reversal program is a separate theorem route.  Modular centered operators, motion/contact rows, nonconvex Coxeter-pure rows, the parked exceptional \(S_6\) benchmark, and the blocked numerical bridge remain future or parked work unless their own gates promote them.

\begin{table}[ht]
\centering
\caption{Future routes and open debts kept outside the present claim boundary.}
\label{tab:future-routes}
\begin{tabular}{@{}L{0.18\textwidth}L{0.32\textwidth}L{0.40\textwidth}@{}}
\toprule
route or debt & current role & admission condition before promotion \\
\midrule
signed-reversal route & separate theorem-level route & must close its own remaining gates; kept pointer-only here \\
P16 modular centered operators & future modular route & needs finite-shadow channel and modular proof boundary \\
future grammar extensions & future grammar branches & need source locks and non-overlap with this paper's claims \\
M19 and \texttt{lat565} & controlled X05 exhibit & mechanism and heterogeneous exceptions must be characterized separately \\
shifted-diagonal row & balance-control exhibit & all-\(n\) statement remains a proof obligation \\
5568 bridge candidate & blocked numerical coincidence/bridge candidate & requires an identified morphism or explicit non-bridge closeout \\
exceptional \(S_6\) benchmark & parked benchmark & no theorem import without its own prior-art and source-lock gates \\
\bottomrule
\end{tabular}
\end{table}

The X05/M19 exhibit has its own boundary.  This paper may use it to show that \(\PhiFCIG\) can recover a recurring hidden group in a centered-operator rank locus.  This paper does not claim a general Sudoku theorem, a full mechanism theorem for the rank drop, or an explanation of the heterogeneous out-of-\(G_0\) exceptions such as \texttt{lat565}.

The shifted-diagonal material and the 5568 bridge candidate remain fenced.  The shifted-diagonal row is useful as a balance-control exhibit, but the all-\(n\) statement is still a proof obligation.  The numerical bridge involving 5568 is not used as a theorem or as evidence of an identified morphism.

This boundary is part of the result of the interface.  The grammar is useful only if it admits source-locked finite rows while refusing rows whose finite shadow, representation channel, or fingerprint package is not yet defined.  If later drafting removes the worked examples in \cref{tab:type-a-rankmass,tab:p12-value-witness,tab:signed-rankmass,tab:p14-value-witness,tab:phi-payoff,tab:m19-census,tab:abd-contrast,tab:future-routes}, the honest fallback is a shorter interface note rather than a full paper.